
This study tested whether lifecycle-stage functional categories from dataset documentation frameworks manifest as unsupervised cluster structure in semantic embeddings across HuggingFace, OpenML, and UCI repositories. Linear probes achieved 79.6\% cross-repository accuracy, confirming that lifecycle signals are encoded in embeddings. However, K-means clustering achieved NMI of 0.023, failing to exceed baseline methods and falling 96\% short of the threshold for meaningful recovery.

The supervised-unsupervised performance gap indicates that lifecycle categories exist as detectable signals but do not form natural clusters under current metadata distributions. Class imbalance (8.3\% Responsible AI fields, 11:1 ratio) and repository heterogeneity (30-fold NMI variance between repositories) prevent standard clustering algorithms from recovering lifecycle structure.

These findings indicate that label-free cross-repository lifecycle detection via unsupervised clustering is not feasible. However, supervised and semi-supervised approaches---few-shot learning, active learning, label propagation---remain viable. Preliminary evidence suggests 60 training samples enable 76\% generalization to held-out repositories, though production deployment requires validation at scale.

The boundary identified---supervised signals exist but do not cluster naturally---clarifies when unsupervised discovery fails (severe class imbalance, repository heterogeneity) and redirects research toward efficient label propagation for ecosystem-scale metadata quality assessment.

